\documentclass[11pt]{article}
\usepackage[a4paper,margin=1in]{geometry}
\usepackage{amsmath,amssymb,mathtools,bm}
\usepackage{enumitem,booktabs,array}
\usepackage{tcolorbox}
\usepackage{hyperref}
\usepackage[protrusion=true,expansion=false]{microtype}
\hypersetup{colorlinks=true,linkcolor=blue,urlcolor=blue,citecolor=blue}
\setlist[itemize]{topsep=2pt,itemsep=2pt}
\setlist[enumerate]{topsep=2pt,itemsep=2pt}
\newtcolorbox{keybox}{colback=blue!3!white,colframe=blue!40!black,boxrule=0.4pt,arc=1mm}
\newtcolorbox{alertbox}{colback=red!3!white,colframe=red!40!black,boxrule=0.4pt,arc=1mm}
\newtcolorbox{answerbox}{colback=green!3!white,colframe=green!40!black,boxrule=0.4pt,arc=1mm}
\newtcolorbox{drillbox}{colback=yellow!5!white,colframe=orange!60!black,boxrule=0.4pt,arc=1mm}
\newcommand{\EV}{\mathbb{E}}
\newcommand{\Var}{\mathrm{Var}}
\newcommand{\Cov}{\mathrm{Cov}}
\newcommand{\blank}[1]{\underline{\hspace{#1}}}

\title{W09-04: Géczy, Minton \& Schrand (2007, JF) \\
\large ``Taking a View: Corporate Speculation, Governance, and Compensation'' --- Active Recall Workbook}
\author{Banking \& Risk Management --- Exam Prep}
\date{\today}
\begin{document}\maketitle

\begin{keybox}
\textbf{Paper in one sentence.} Based on a 1998 survey of 186 Fortune 500 derivatives users, GMS show that roughly half of firms that use derivatives \emph{speculate} (take ``views''), and that such firms are distinguished by \emph{stronger governance}, \emph{tighter internal controls}, and \emph{performance-linked compensation} --- so speculation is disciplined where managers have the tools to understand and measure the risk.

\textbf{Placement.} Completes the hedging-vs-speculation set: BCH (2006, W09-03) shows selective hedging does not add value; GMS (2007) shows \emph{who} speculates and \emph{under what governance conditions}. Complement to GMS (1997, W07-01) on the determinants of derivative use.
\end{keybox}

\section*{Round 1: Complete Walk-Through}

\subsection*{1.1 Data}
1998 survey of 186 non-financial Fortune 500 derivatives users, augmented with compensation, governance, and fundamentals data. Key survey items: do you ``frequently'' or ``sometimes'' take a view; how do you measure derivative performance; what governance controls are in place.

\subsection*{1.2 Measures of speculation}
\begin{itemize}
\item Self-reported frequency of ``taking a view''.
\item Whether derivative performance is separately evaluated.
\item Whether compensation is linked to derivative outcomes.
\end{itemize}
About 50\% of derivative users report some speculation; $\sim 15\%$ report ``frequent'' view-taking.

\subsection*{1.3 Determinants of speculation}
Logit of speculation dummy on governance and compensation measures:
\begin{itemize}
\item Firms with more sophisticated treasury functions speculate more.
\item Firms with tighter internal controls (audit, board oversight) speculate more.
\item Firms with pay linked to hedging/derivative outcomes speculate more.
\item Firms with stronger governance (board independence, blockholders) speculate more.
\end{itemize}

\subsection*{1.4 Interpretation}
The headline is counterintuitive: speculation is not the behaviour of cowboy firms, but of firms that have the governance and control apparatus to \emph{discipline} speculative activity. Read with BCH: even when disciplined, speculation does not generate measurable value, but the governance ensures it does not destroy value either.

\subsection*{1.5 Caveats}
\begin{alertbox}
\begin{itemize}
\item Self-report in a survey: possible strategic answers.
\item Non-response bias: firms that refuse the survey may differ.
\item Cross-section: cannot establish causality between governance and speculation.
\end{itemize}
\end{alertbox}

\section*{Round 2: Level 1 Fill-in}
\begin{enumerate}
\item Survey year: \blank{1cm}; sample size: \blank{1cm}.
\item Fraction of users that ``take a view'' sometimes: $\sim$\blank{1cm}\%.
\item Speculating firms have \blank{1cm} governance.
\item Speculating firms often have pay \blank{1cm} to derivative outcomes.
\item Complement paper on whether speculation adds value: Brown--Crabb--\blank{2cm} (2006).
\end{enumerate}

\section*{Round 3: Level 2 Prompts}
\begin{enumerate}
\item Write a stylised logit for speculation and state the expected signs.
\item Explain why stronger governance enables rather than prevents speculation.
\item Contrast GMS (2007) and Brown--Crabb--Haushalter (2006).
\item How does GMS change our view of ``taking a view''?
\end{enumerate}

\section*{Round 4: Minimal Scaffolding}
From a blank page: (i) survey \& sample, (ii) speculation prevalence, (iii) determinants, (iv) governance interpretation, (v) two caveats.

\section*{Round 5: Blank Page}
Essay: the role of governance in corporate speculation --- GMS's counterintuitive story.

\vspace{6pt}
\begin{drillbox}
\textbf{Speed Drill.} (1) Survey year/size. (2) Fraction speculating. (3) Governance relationship. (4) Compensation link. (5) Complement paper.
\end{drillbox}

\section*{Quick Reference \& Answer Key}
\begin{answerbox}
\begin{itemize}
\item \textbf{Survey.} 1998, 186 Fortune 500 derivative users.
\item \textbf{Prevalence.} $\sim$50\% speculate sometimes; $\sim$15\% frequently.
\item \textbf{Determinants.} Stronger governance $+$ comp tied to derivative P\&L.
\item \textbf{Interpretation.} Speculation is disciplined where controls exist.
\end{itemize}
\end{answerbox}

\begin{keybox}
\textbf{30-second exam pitch.} Géczy, Minton \& Schrand (2007) survey 186 Fortune 500 non-financial derivative users in 1998 and discover that about half sometimes take a ``view'' with their derivatives --- i.e., speculate. The unexpected finding is that speculating firms have \emph{stronger}, not weaker, governance: more sophisticated treasury functions, tighter internal controls, board oversight of derivatives, and compensation tied to derivative outcomes. The interpretation is that speculation appears where firms have built the infrastructure to discipline it. Read alongside Brown--Crabb--Haushalter (2006), which shows selective hedging does not add value in gold miners, GMS suggests a sobering equilibrium: firms that speculate have the governance to prevent losses, and in the cross-section the mean P\&L is near zero --- speculation is permitted where it is carefully controlled, not where it is skilled.
\end{keybox}

\end{document}
